\documentclass[11pt]{article}
\pdfinfoomitdate=1
\pdftrailerid{}
\pdfsuppressptexinfo=15
\usepackage[margin=0.9in]{geometry}
\usepackage{amsmath,amssymb,amsthm,graphicx}
\usepackage[T1]{fontenc}
\usepackage{lmodern}
\usepackage[title]{appendix}
\usepackage[colorlinks=true,linkcolor=blue,citecolor=blue,urlcolor=blue]{hyperref}
\hypersetup{pdftitle={A counterexample to the cocircuit diameter conjecture},
 pdfauthor={P. Graham},pdfsubject={Cocircuit graphs of realizable uniform oriented matroids},
 pdfkeywords={oriented matroid, cocircuit graph, diameter, uniform perturbation}}
\newtheorem{theorem}{Theorem}
\newtheorem{lemma}[theorem]{Lemma}
\theoremstyle{remark}
\newtheorem{remark}[theorem]{Remark}
\newcommand{\Q}{\mathbb Q}
\newcommand{\R}{\mathbb R}
\newcommand{\diam}{\operatorname{diam}}
\newcommand{\sgn}{\operatorname{sgn}}
\title{A counterexample to the cocircuit diameter conjecture}
\author{P. Graham\\ \small\url{https://orcid.org/0009-0007-8493-1505}}
\date{October 26, 2026}
\begin{document}
\maketitle
\begin{abstract}
A rationally realizable uniform rank-41 oriented matroid on 80 elements
has cocircuit-graph diameter at least $42>41=n-r+2$, refuting the
general cocircuit diameter conjecture. We build on the rank-21 example
of Adler, De Loera, Klee and Zhang, the non-Hirsch polytope of
Matschke, Santos and Weibel, Santos's product construction, and the
source certification of Allamigeon, Canu and Strub. A count of negative
coordinates confines every potentially shorter walk to two affine sections;
projection to the product factors then gives the lower bound.
Two rational perturbations make the resulting oriented matroid uniform.
Lean~4 separately verifies the matrix existence and walk lower bound
for a different realizing matrix; the printed perturbations are proved here,
not checked in Lean.

This work has not yet been reviewed by experts.
The proofs, the Lean formalization, the certificate checkers and the text were produced by AI models
(from OpenAI and Anthropic), working under the author's direction.
\end{abstract}
\medskip

\noindent 2020 \emph{Mathematics Subject Classification}: 52C40, 05C12, 52B05.\\
\emph{Keywords}: cocircuit graph, oriented matroid, diameter, uniform perturbation.

\section{Introduction}

Let $A\in\R^{n\times r}$ have rank $r$. Its sign covectors are
$\sgn(Ax)$, $x\in\R^r$, with values in $\{-1,0,1\}^n$.
Its \emph{cocircuits} are the nonzero covectors of minimal support.
Write $X^0=\{i:X_i=0\}$ and
$S(X,Y)=\{i:X_i=-Y_i\ne0\}$ for the separation set. When every $r$-row minor of $A$ is nonzero,
the associated oriented matroid is \emph{uniform}; every cocircuit has
$r-1$ zeros. Its \emph{full cocircuit graph} $G^*(A)$ has all signed cocircuits
as vertices, with distinct $X,Y$ adjacent when
$S(X,Y)=\varnothing$ and $|X^0\cap Y^0|=r-2$.
Rational realizability means $A$ can be chosen over $\Q$.
Cocircuits $X,Y$ are \emph{antipodal} if $Y=-X$, and nonantipodal otherwise.
Write $d_H$ for distance in $H$.

Adler, De Loera, Klee and Zhang's (ADKZ) Conjecture~1 states that every rank-$r$ oriented matroid on
$n$ elements has cocircuit-graph diameter at most $n-r+2$
\cite[Conjecture~1]{adkz}. Babson, Finschi and Fukuda's early
farthest-pairs question \cite[final remarks, pp.~13--14]{bff} and
Felsner, G\'omez, Knauer, Montellano-Ballesteros and Strausz's
Question~4.2 \cite{felsner} precede ADKZ's formulation. ADKZ report Fukuda's personal communication
that the bound is folklore. Klee's report attributes its introduction
to Fukuda and Terlaky \cite[p.~2425]{klee}.

\begin{theorem}\label{thm:main}
There exists a rationally realizable uniform oriented matroid of rank 41
on 80 elements whose cocircuit graph has diameter at least
$42>41=80-41+2$.
\end{theorem}

ADKZ state a uniform rank-21 example on 40 elements with nonantipodal
cocircuits at distance at least 21 \cite[Proposition~8]{adkz}.
Their proof uses the 20-dimensional non-Hirsch polytope of Matschke,
Santos and Weibel (MSW), with 40 facets
\cite[Theorem~1.3]{msw}, obtained using Santos's construction
\cite[Theorem~2.6]{santos}. Allamigeon, Canu and Strub (ACS) proved
its diameter 21 in Coq \cite[\S5, PDF p.~12]{acs}.
We reuse their numerical inequality data and check a new certificate for
the designated pair. Product amplification is standard
\cite[Lemma~6.1 and \S6]{santos}, where Santos cites Klee--Walkup,
Proposition~1.3.

We must control walks that leave the product's nonnegative section.
The lifted source matrix $M_0$ is not uniform: rows 5, 7--10 and 27--29
are supported on the first six columns (Appendix~\ref{app:certificate}).
We first perturb it to a uniform rank-21 factor, preserving
the pair at distance 21. Two factors give 80 inequalities in rank 41:
the product pair is 42 edges apart; the conjectured bound is
$80-41+2=41$. To exclude shorter cocircuit walks,
Lemma~\ref{lem:count} generalizes ADKZ's count of sign-coordinate changes
\cite[Lemma~5 and proof of Proposition~8]{adkz}: a walk of at most 41
edges can visit at most one negative coordinate. Such a walk stays in
the union of two sections. A second uniform perturbation preserves
these sections; projection to the factors still requires $21+21$ edges.

\section{A lower bound on walk length}

Fix an inequality matrix $M\in\R^{n\times r}$ and put
$C_0(M)=\{z:Mz\ge0\}$ and
$C_e(M)=\{z:(Mz)_e\le0,\ (Mz)_i\ge0\ (i\ne e)\}$.
Choose a rational linear functional $c$ positive on every nonzero vector
in these cones. The \emph{affine section} of a cone $C$ is $C\cap\{c(z)=1\}$;
positivity makes it bounded (Lemma~\ref{lem:stability}(a)).
The sections of $C_0$ and $C_e$ are
nonnegative and singly reversed, respectively. A row is \emph{tight} when equality holds;
the tight row labels form the \emph{active set}.
When $M$ is uniform and the section is full-dimensional, we identify
a section vertex $z$ with the cocircuit $\sgn(Mz)$; its active set is
its zero set (Lemma~\ref{lem:cells}).
A $d$-dimensional section is
\emph{simple} if each vertex lies on exactly $d$ facets.

\begin{samepage}
\begin{lemma}\label{lem:count}
Let $G^*$ be the cocircuit graph of a uniform rank-$(d+1)$ oriented
matroid on $2d$ elements. Suppose that $X^0,Y^0$ are complementary
$d$-element sets and $X,Y$ are positive elsewhere. For a walk of length
$L$ from $X$ to $Y$, let $k$ be the number of coordinates that are
negative at any vertex of the walk. Then
\[
                      2d+2k\le2L.
\]
\end{lemma}
\end{samepage}
\begin{proof}
Adjacent cocircuits have $d$ zeros each, sharing $d-1$.
Exactly one zero becomes nonzero and one nonzero becomes zero.
All other nonzero entries agree because the separation set is empty.
Thus exactly two coordinates change, each by absolute amount one.
The complementary endpoint zero sets make every coordinate go from
$0$ to $1$ or from $1$ to $0$. Its total absolute change is at least one,
or at least three if it visits $-1$. Summing gives $2d+2k\le2L$.
\end{proof}

For $d=40$ and $L\le41$, the walk uses at most one fixed negative
coordinate $e$, so every vertex lies in $C_0$ or $C_e$.
If none is negative, any $e$ works.

\begin{lemma}\label{lem:cells}
Let $M\in\Q^{n\times r}$ have rank $r\ge2$ and be uniform. Let a rational
linear functional $c$ be positive on every nonzero vector of $C_0(M)$
and every $C_e(M)$. Suppose each section by $c(z)=1$ is bounded
and has dimension $r-1$.
Each section graph is the restriction of $G^*(M)$ to cocircuits satisfying its
sign requirements $X_i\ge0$ for $C_0$, or $X_e\le0$ and $X_i\ge0$
for $i\ne e$ for $C_e$. For every row index $e$, every walk whose negative coordinates
are contained in $\{e\}$ lies in the union of the graphs of the sections
of $C_0(M)$ and $C_e(M)$.
\end{lemma}
\begin{proof}
The $r-1$ zero rows of a cocircuit are independent. Their kernel is a line;
choosing the signs selects a ray. Positivity of $c$ gives exactly one
representative with $c(z)=1$, and excludes the antipodal ray from that cone.
This representative is a section vertex.
Conversely, a section vertex needs $r-1$ independent
tight rows; uniformity forbids more, so its sign vector is a cocircuit.
These $r-1$ independent active rows give local coordinates in the
$(r-1)$-dimensional section, so it is simple. Two vertices sharing $r-2$ tight rows
are the endpoints of their common one-dimensional face, and every edge
has such endpoints. Conversely, full-graph adjacency between section
vertices requires these $r-2$ shared zeros, so gives such a section edge.
Separation is empty inside one of these cones; thus its
graph is the induced cocircuit subgraph on those vertices.
For an adjacent pair with no negative coordinates outside $e$, if either
has a negative entry at $e$, the other has a nonpositive entry there:
opposite nonzero signs are excluded by adjacency. Both then lie in $C_e$.
Otherwise both lie in $C_0$. When both signs at $e$ vanish, both vertices lie on the shared facet
and the edge lies in both graphs. Every walk edge therefore belongs to
one of the two section graphs.
\end{proof}

\section{The construction}

Put $P=C_0(M_0)\cap\{z:c(z)=1\}$, with the rational matrix $M_0$ and
functional $c$ specified in Appendix~\ref{app:certificate}.
The source certificate \cite{supplement}, detailed in Appendix~\ref{app:cover},
gives a simple 20-dimensional $P$ and $d_P(u,v)=21$ for vertices $u,v$
with active rows 1--20 and 21--40, respectively.
The inequality data are \texttt{poly20dim21.ine} in Coq-Polyhedra
(Allamigeon, Katz and Strub). ACS attribute the supplied inequalities
to Weibel \cite[\S5, PDF p.~11, footnote~6]{acs}; Allamigeon, Id-Sahra
and Strub likewise credit MSW for the 20- and 23-dimensional non-Hirsch
polytopes they consider \cite[\S1]{cert26}.
We use the documented attribution for credit; the source file was not compared
entrywise with MSW's table. The proof uses the exact matrix and certificate
specified here; the supplement records the source version and comparisons.
Clearing each row's reduced denominators by their least common multiple gives
$M_0\in\mathbb Z^{40\times21}$, leaving every sign unchanged.
Let
$\sigma_e\in\{-1,1\}^{40}$ have entries $\sigma_{e,i}$:
$\sigma_0$ is all positive, and $\sigma_e$ reverses row $e$ for $e>0$.
Appendix~\ref{app:certificate} gives a rational row vector $c$ and weights
$\lambda_{e,i}\ge1$ satisfying
\[
 c=\sum_{i=1}^{40}\lambda_{e,i}\sigma_{e,i}(M_0)_i,
 \qquad(0\le e\le40).
\]
These exact identities, with $\operatorname{rank}M_0=21$, imply
$c(z)>0$ on every nonzero vector of each closed cone: its signed slacks
are nonnegative and not all zero.


\begin{lemma}[Stability of bounded simple sections]\label{lem:stability}
Let $M\in\R^{n\times r}$ and let $\Sigma\subseteq\{-1,1\}^n$ be finite.
Suppose a linear functional $c$ is positive on every nonzero vector of
\[
 C_\sigma(M)=\{z:\sigma_i(Mz)_i\ge0\text{ for all }i\},
 \qquad \sigma\in\Sigma.
\]
\textup{(a)} For all $M'$ in a sufficiently small neighborhood of $M$ in
$\R^{n\times r}$, the same positivity
holds and the sections $C_\sigma(M')\cap\{c=1\}$ are uniformly bounded.
\textup{(b)} For each $\sigma$ whose original section has dimension $r-1$ and
exactly $r-1$ independent active rows at every vertex, its perturbed
section has the same dimension and face lattice,
with the same row labels and nonzero signs. For two patterns
$\sigma,\sigma'$ satisfying these hypotheses, their perturbed sections
intersect in the face where all rows on which $\sigma$ and $\sigma'$
differ are tight. This is the continuation of their original common face,
and the two continuation maps agree on its vertices and edges.
\end{lemma}
\begin{proof}
\emph{Positivity and boundedness.} Failure of positivity along $M'\to M$ would give unit vectors in
perturbed cones with $c\le0$; a convergent subsequence contradicts the
original positivity. Unbounded normalized sections similarly give a unit
limit in an original cone with $c=0$.
\emph{Vertices.} For a full-dimensional original section, a point with all inequalities
strict remains feasible, preserving dimension. Each original vertex continues uniquely through its nonsingular
active-row system and $c=1$. Every inactive signed slack was strictly
positive; by continuity it remains so, giving exactly the original active set.

Suppose additional vertices appeared along $M'\to M$. Uniform boundedness
gives a convergent subsequence with limit $z$. Finiteness of row
subsets gives a further subsequence with a fixed set $I$ of $r-1$
independent tight rows at the perturbed vertices. These rows need not
a priori be independent at $M$. The original bounded face with $I$
tight contains $z$ by continuity of every signed inequality and tight
equation. This nonempty bounded polyhedral face has a vertex $w$.
Its active set contains $I$ and has the same cardinality $r-1$, so equals $I$. Thus $M_I$ has rank
$r-1$, and $M_Iz=0$, $c(z)=1$ uniquely determines $w$.
For sufficiently small perturbations the system remains nonsingular;
the alleged additional vertex is the continuation of $w$.

\emph{Faces.} Identify each vertex with its unique continuation, retaining its active
set and nonzero signs. For a row subset $I_0$, the face cut out by making
$I_0$ tight has precisely the vertices whose active sets contain $I_0$.
The continuation map preserves this vertex set and its emptiness,
for each section satisfying the stated dimension and active-row hypotheses.
Every face of an inequality polytope has this form and is the convex hull
of its vertices. Thus equal faces, and inclusion between faces, are
preserved in both directions. This gives a face-lattice bijection,
preserving edges and their endpoint pairs.

\emph{Common faces.} If two patterns differ on $I_0$, the intersection of their sections is
the face obtained by making $I_0$ tight in either section. Every vertex of this face is a vertex of both sections: in a convex
decomposition into section points, all signed slacks on $I_0$ are nonnegative
and average to zero, so both points lie in the common face. A nontrivial
decomposition would contradict its being a vertex there.
A row is tight precisely
when its linear form is zero, independently of the sign pattern.
The active set $I$ is thus the same in both sections, has exactly
$r-1$ rows, and contains $I_0$. Its continuation is the unique
solution of $M'_Iz=0$, $c(z)=1$ in either section. The identifications,
full sign vectors and common-face edges consequently agree.
Intersecting the neighborhoods for the finitely many sections and vertex
systems gives one neighborhood for every $\sigma\in\Sigma$.
\end{proof}

Set $F=M_0+\varepsilon J$, where $J_{ij}=(i+1)^{j-1}$ has 21 columns.
Each 21-row minor is a polynomial in $\varepsilon$ with nonzero
Vandermonde leading coefficient, so all are nonzero for sufficiently
small rational $\varepsilon>0$. Apply Lemma~\ref{lem:stability} with
$\Sigma=\{\sigma_0,\sigma_1,\ldots,\sigma_{40}\}$: $c$ stays positive
on these 41 cones.
Each vertex of $P$ has exactly 20 independent tight rows, as the
certificate checks. The section $P$ of $C_0(M_0)$ is therefore simple, so
Lemma~\ref{lem:stability} preserves its graph with the active row labels.
Let $F\in\Q^{40\times21}$ be the resulting uniform matrix, and $U,V$
the unique continuations of $u,v$. Write $H_0,H_e$ for the section
graphs of $C_0(F),C_e(F)$, respectively.
\begin{lemma}\label{lem:facetcross}
Each reversed section $C_e(F)\cap\{c=1\}$ has dimension 20,
and every vertex has exactly 20 independent tight rows.
\end{lemma}
\begin{proof}
All reversed sections are bounded by Lemma~\ref{lem:stability}, applied
above with all 41 patterns. (1) Every row is active at $U$ or $V$. There the twenty independent
active constraints are local coordinates on $c=1$. Keeping one zero
and making the other nineteen positive gives a 19-dimensional face.
Thus every row defines a facet of the nonnegative section.

(2) A facet's affine span in $c=1$ spans a linear hyperplane in
$\R^{21}$. Two rows defining the same facet therefore vanish on the same
hyperplane and are proportional, contradicting a nonzero 21-row minor
of the uniform matrix $F$. Distinct rows define distinct facets.

(3) At a relative-interior point of facet $e$, any other tight row would
vanish on the facet's affine span: a nonnegative affine form vanishing
at an interior point vanishes identically there. This would make its
facet equal to facet $e$, since both are facets by (1), contradicting (2). Thus only row $e$ is tight.
Crossing its hyperplane within $c=1$ makes row $e$ negative and all
others positive, giving a 20-dimensional reversed section.

(4) At every reversed-section vertex at least twenty independent
constraints are tight. Uniformity forbids 21 tight rows, whose common
kernel would be zero. Thus exactly twenty rows are tight and independent.
\end{proof}
By Lemma~\ref{lem:cells}, $H_0,H_e$ are subgraphs of $G^*(F)$,
identified by sign vectors. Their union is the induced subgraph on the
cocircuits satisfying either set of sign inequalities.
$H_0$ inherits $P$'s graph and row labels,
giving $d_{H_0}(U,V)=21$.
Their active labels still partition the 40 facets, so the cocircuits
corresponding to $U,V$ have complementary zero sets and are positive elsewhere.

\begin{lemma}\label{lem:factor}
Every $U$--$V$ walk in $G^*(F)$ has length at least 21.
\end{lemma}
\begin{proof}
This is ADKZ's Proposition~8 count applied to $F$.
For a walk of length at most 20, Lemma~\ref{lem:count} with $d=20$
gives $k=0$. Its vertices are nonnegative, so Lemma~\ref{lem:cells} puts it
in $H_0$, contradicting $d_{H_0}(U,V)=21$.

\end{proof}

Choose rational coordinates $(s,x)\in\R\times\R^{20}$ with $s=c(z)$,
so the affine sections lie in $s=1$. After this rational invertible
change of coordinates we may assume that
$F=(\beta\ R)$ with $\beta\in\Q^{40}$ and $R\in\Q^{40\times20}$,
so $F(s,x)=\beta s+Rx$. The change preserves all sign vectors and
nonvanishing maximal minors. In product coordinates $(s,x,y)$,
the inequality matrix is
\[
 B=\begin{pmatrix}\beta&R&0\\\beta&0&R\end{pmatrix}\in\Q^{80\times41}.
\]
The section of $C_0(B)$ is a product of two 20-dimensional polytopes.
Its vertices are pairs of factor sign vectors; each edge moves one factor.
For a flip in the first factor, the reversed section graph is
$H_e\square H_0$. We use only the inclusion
\[
 (H_0\square H_0)\cup(H_e\square H_0)
 \ \subseteq\ (H_0\cup H_e)\square H_0.
\]
Here $\square$ is the Cartesian product of graphs. The second-factor case
is symmetric.


For the second perturbation, set $W\in\Q^{80\times41}$ with
$W_{ij}=(i+1)^{j-1}$ and seek
$A=B+\delta W$ with rational $\delta>0$ sufficiently small.

\begin{lemma}\label{lem:preserve}
For some rational $\delta>0$, the matrix $A=B+\delta W$ is uniform.
The functional $s$ is positive on every nonzero vector of $C_0(A)$ and
all $C_e(A)$. Their sections by $s=1$ are bounded, have dimension 40,
and have the same
graphs and common faces as the corresponding sections for $B$,
identifying vertices by active rows together with their sign vectors.
\end{lemma}
\begin{proof}
Fix $F$ and $s$ before choosing $\delta$. Every factor section has dimension 20,
by the construction of $P$ for $H_0$ and Lemma~\ref{lem:facetcross} for $H_e$.
Uniformity of $F$ gives exactly 20 independent active rows at each vertex,
with all other signed inequalities strict. All 81 sections (the nonnegative section and one reversal for each
of the 80 rows) thus have dimension 40.
Fixing $s=1$ makes the two factor blocks independent. Each product-section
vertex is a pair of factor-section vertices. The active equations in
the $x$ and $y$ variables form two nonsingular blocks of size 20,
giving exactly 40 independent active equations. The remaining signed
inequalities are strict.
On each product cone the two factor vectors belong to $C_0(F)$ or one
$C_e(F)$ and share the coordinate $s$. Positivity of the factor functional
forces $s\ge0$; if $s=0$, both factor vectors vanish. Thus $s>0$ on
every nonzero vector of all 81 product cones. Their sections in the
common hyperplane $s=1$ are bounded. The product may be nonuniform.
For any 41-row set $I$,
$\det(B_I+\delta W_I)$ has degree 41: its top coefficient $\det W_I$ is the
Vandermonde determinant of 41 distinct integers and is nonzero.
All finitely many maximal-minor polynomials are therefore nonzero for
sufficiently small $\delta>0$.
Use $\Sigma$ consisting of the all-positive pattern and each of its
80 single reversals. Lemma~\ref{lem:stability} applies: the functional $s$ remains positive, the sections remain
uniformly bounded, and their labeled graphs and shared faces are preserved.
\end{proof}


\begin{proof}[Proof of Theorem~\ref{thm:main}]
Lemma~\ref{lem:preserve} gives a uniform $A$, positivity of $s$, and
bounded sections of dimension 40. Thus Lemma~\ref{lem:cells} applies.
In the section of $C_0(A)$ choose the vertices labeled $(U,U)$ and
$(V,V)$. Their zero sets are complementary 40-element sets and
both are positive elsewhere. Suppose
there were a full-graph walk of length $L\le41$ between them.
Lemmas~\ref{lem:count} and~\ref{lem:cells} confine its vertices and edges
to the sections of $C_0(A)$ and one fixed $C_e(A)$.
Lemma~\ref{lem:preserve} transports the walk to $B$, preserving active
rows and sign vectors. The common-face clause of
Lemma~\ref{lem:stability}(b) makes the two section maps agree on every
shared vertex and edge.
If the fixed coordinate $e$ belongs to the first factor,
the resulting walk is in $(H_0\cup H_e)\square H_0$.
The first-factor projection lies in $H_0\cup H_e$ and the second in
$H_0$. By Lemma~\ref{lem:cells} applied to $F$, both $H_0$ and $H_e$
are subgraphs of $G^*(F)$.
Project the whole walk to each factor and delete consecutive repetitions. Each
edge moves exactly one factor, so the lengths of the projected walks
sum to $L$. Lemma~\ref{lem:factor}
requires at least 21 first-factor edges, and
$d_{H_0}(U,V)=21$ requires at least 21 second-factor edges. Thus $L\ge42$, a contradiction. The second-factor case is identical.
Every 41-row minor of $A$ is nonzero, so its rank is 41. Its real sign
covectors give the standard realizable oriented matroid, uniform by these
minors and rationally realizable since $A$ is rational.
\end{proof}

\begin{remark}\label{rem:path}
The distance between these two vertices is exactly 42. The two
21-edge paths in $P$, continued to $H_0$, give the path
\[
 (U,U),\ldots,(V,U),\ldots,(V,V)
\]
in the product section. Lemma~\ref{lem:preserve} preserves its edges
after the second perturbation. This gives the upper bound 42, matching
the lower bound in the proof of Theorem~\ref{thm:main}.
\end{remark}





\section{Formal verification}

The Lean theorem
\nolinkurl{ADKZPricing.Generated.source_frozen_endpoint}
asserts that there is
$A'\in\Q^{80\times41}$ with every 41-row minor nonzero and real signed
cocircuits $X,Y$ having complementary 40-element zero sets and positive
entries elsewhere, such that no walk of at most 41 edges joins them.
This is Theorem~\ref{thm:main} for a different matrix from Section~3.
The walks allow arbitrary signed intermediate cocircuits and repetitions;
antipodal sign vectors are distinct vertices.

In the formal definitions, cocircuits are nonzero covectors minimal
under conformality: a conformal face $Z$ of $X$ has
$Z_i\in\{0,X_i\}$ for every $i$. Adjacency is defined by composition,
$(X\circ Y)_i=X_i$ if $X_i\ne0$ and $Y_i$ otherwise: for distinct
$X,Y$, the compositions $X\circ Y$ and $Y\circ X$ must agree and be a
covector having exactly $X,Y$ as cocircuit conformal faces. In the uniform
case, Lean proves that conformality-minimality agrees with support-minimality
and that this adjacency agrees with Section~1: empty separation set and
39 common zeros. Thus the formal walk bound has the stated full-graph
meaning; the passage from real matrix sign covectors to the standard
realizable oriented matroid is the usual mathematical interpretation.

Lean uses a different positive functional for the same $M_0$ and the
same product-and-perturb strategy. Its rational perturbations are chosen
nonconstructively; the endpoint theorem specifies $A'$ existentially,
without printing its entries.
The printed Vandermonde perturbations are not checked in Lean; their
proofs and Remark~\ref{rem:path}'s upper bound are written. The formal
statement, axioms and reproduction instructions are given in
\cite{supplement}.

\section{Further questions}

The diameter of the full cocircuit graph of this example is unknown. What is the least counterexample rank?
What sharp bound replaces $n-r+2$?

Fukuda's linear-bound question remains \cite[\S3.7.2, PDF p.~16]{survey};
see also \cite[Open Problem~5]{finschi}.

\begin{appendices}
\section{An exact positive-functional certificate}\label{app:certificate}
% BEGIN EXPLICIT CERTIFICATE
This appendix specifies the source matrix $M_0$ and a positive functional
for the cones obtained by reversing at most one row. Row and column labels
are one-based: $1\le i\le40$, $1\le j\le21$; $e=0$ denotes the
nonnegative cone and $e=1,\ldots,40$ reverses row $e$.

\subsection*{The integer source matrix}
Let $u_j$ denote the $j$th coordinate row of $\mathbb Z^{21}$ and let $m_i$
be row $i$ of $M_0$. For the rows in the following table set
\[
 m_i=s_i(u_1+100\eta_i u_2)+a_i u_3+b_i u_4+p_i u_5+q_i u_6+T_i,
 \qquad \eta_i=\begin{cases}1&i\le20,\\-1&i\ge21.\end{cases}
\]
\begin{center}\small
\setlength{\tabcolsep}{7pt}
\begin{tabular}{r|rrrrr|l}
$i$&$s_i$&$a_i$&$b_i$&$p_i$&$q_i$&$T_i$\\\hline
1&1&0&0&21&-7&$u_{17}+u_{18}$\\
2&1&0&0&16&-15&$u_{18}+u_{19}$\\
3&1&0&0&0&-32&$u_{19}+u_{20}$\\
4&1&0&0&-16&-15&$u_{20}+u_{21}$\\
5&1&0&0&-21&-7&$0$\\
6&1&0&0&-20&-4&$u_{21}$\\
7&1&0&0&0&32&$0$\\
8&1&0&0&20&-4&$0$\\
9&100&3&-2&0&-3000&$0$\\
10&100&-3&-2&0&3000&$0$\\
11&2000&-3&7&0&63600&$2000u_{16}$\\
21&1&30&0&0&0&$u_{9}+u_{10}$\\
22&1&4&-15&0&0&$u_{10}+u_{11}$\\
23&2&0&-33&0&0&$2u_{11}+2u_{12}$\\
24&1&-1&-16&0&0&$u_{12}+u_{13}$\\
25&2&-55&0&0&0&$2u_{13}+2u_{14}$\\
26&1&-17&18&0&0&$u_{14}$\\
27&1&0&38&0&0&$0$\\
28&1&22&17&0&0&$0$\\
29&5&-50&0&1&-1&$0$\\
30&100&2999&0&-12&-20&$100u_{9}$\\
31&10000&299999&0&0&100&$10000u_{8}$\\
32&20000&-549000&0&4&25&$20000u_{7}$\\
\end{tabular}
\end{center}
The remaining seventeen rows are specified by two triangular recipes. Define
\[
 \begin{aligned}
 \alpha&=2000u_1+200000u_2+3u_3+7u_4-63600u_6+2000u_{15},\\
 (t_7,\ldots,t_{14})&=2(10^{10},10^{10},10^{10},10^{13},10^{14},10^{14},10^{14},10^{14}).
 \end{aligned}
\]
Then
\[
 m_{12}=\alpha+\sum_{k=7}^{14}t_ku_k,\qquad
 m_{j+6}=\alpha+\sum_{k=7}^{j-1}t_ku_k-t_ju_j\quad(7\le j\le14).
\]
Likewise define
\[
 \begin{aligned}
 \gamma&=2000u_1-200000u_2-54000u_3+4u_5-25u_6,\\
 (h_{15},\ldots,h_{21})&=2(10^8,10^{10},10^{10},10^{10},10^{11},10^{11},10^{12}).
 \end{aligned}
\]
Then
\[
 m_{33}=\gamma+\sum_{k=15}^{21}h_ku_k,\qquad
 m_{j+19}=\gamma+\sum_{k=15}^{j-1}h_ku_k-h_ju_j\quad(15\le j\le21).
\]
Empty sums are zero, and all entries not specified by these recipes are zero.
These recipes give exactly the $40\times21$ integer matrix obtained by clearing
each source row's reduced denominators by their least common multiple.
\subsection*{The functional and weights}
Set $c(z)=40000\sum_{j=1}^{21}v_jz_j$, where the integers $v_j$ are
\begin{center}\small
\setlength{\tabcolsep}{5pt}
\begin{tabular}{r r|r r|r r}
$j$&$v_j$&$j$&$v_j$&$j$&$v_j$\\\hline
1&$7\,889\,496\,505\,278\,816$&8&$16\,770\,653\,889\,890$&15&$11\,008\,963\,417$\\
2&-$788\,643\,250\,921\,389\,998$&9&$16\,773\,144\,135\,505$&16&$193\,944\,140\,058$\\
3&-322873814&10&$2\,918\,818\,204\,673\,914$&17&$570\,600\,902\,642$\\
4&$73\,889\,176\,985\,687\,161$&11&$2\,039\,103\,229\,854\,373$&18&$574\,335\,730\,279$\\
5&0&12&$1\,964\,406\,684\,567\,592$&19&$186\,741\,363\,964$\\
6&-$8\,399\,586\,269\,177$&13&$2\,477\,284\,371\,671\,404$&20&$112\,044\,818\,677$\\
7&$16\,778\,123\,544\,418$&14&$2\,402\,587\,826\,384\,623$&21&$373\,482\,727\,181$\\
\end{tabular}
\end{center}
Let $\sigma_{0,i}=1$ for every $i$; for $1\le e\le40$ let
$\sigma_{e,e}=-1$ and $\sigma_{e,i}=1$ for $i\ne e$.
For each $e$, specify an increasing list $B_e=(\ell_1,\ldots,\ell_{21})$ as follows:
\[
 \begin{gathered}
 K=\{2,5,7,13,14,15,16,17,18,19,20,21,25,27,33,35,36,37,38,39,40\},\\
 B_e=\operatorname{sort}\bigl((K\setminus R_e)\cup S_e\bigr).
 \end{gathered}
\]
In the table, $\mathcal D=\{2,4,6,8,10,11,22,23,24,26,28,29,30,31,32\}$.
Each displayed comma-separated entry denotes a set, and a multi-label row
assigns the same sets to every listed $e$.
\begin{center}\scriptsize
\setlength{\tabcolsep}{14pt}
\renewcommand{\arraystretch}{1.03}
\begin{tabular}{r|l|l}
$e$&$R_e$ (delete from $K$)&$S_e$ (add to $K$)\\\hline
$0,1$&$\varnothing$&$\varnothing$\\
$\mathcal D$&$2$&$1$\\
$3$&$2,19,25$&$8,12,34$\\
$5$&$2,5,7,38$&$1,6,12,34$\\
$7$&$2,7$&$1,11$\\
$9$&$5,7,33,38,39$&$1,4,6,11,12$\\
$12$&$2,5,27$&$3,4,26$\\
$13$&$2,5,13,16,21,25$&$4,6,12,28,30,34$\\
$14$&$2,14,25,27$&$3,12,28,32$\\
$15$&$2,5,15,21,25,27$&$4,12,26,28,31,34$\\
$16$&$2,5,16,21,25$&$3,12,30,31,34$\\
$17$&$2,5,17,18$&$3,6,22,26$\\
$18$&$2,5,13,18,19$&$3,4,12,22,26$\\
$19$&$2,5,19,25,27$&$3,6,12,26,34$\\
$20$&$2,5,16,18,20$&$3,4,12,22,24$\\
$21$&$2,5,13,15,21,25$&$1,6,12,22,32,34$\\
$25$&$5,7,25,33,38,39$&$1,4,6,12,32,34$\\
$27$&$2,5,19,27$&$1,6,26,34$\\
$33$&$2,5,7,25,33,37$&$1,4,8,11,12,24$\\
$34,38,39$&$5,33,38,39$&$1,4,6,12$\\
$35$&$2,5,7,33,35,38$&$1,3,6,8,11,12$\\
$36$&$5,36,39$&$1,6,12$\\
$37$&$5,36,37$&$3,4,12$\\
$40$&$2,35,38,40$&$1,3,6,12$\\
\end{tabular}
\end{center}
Write $c^{\mathsf T}=40000(v_1,\ldots,v_{21})^{\mathsf T}$ and define
\[
 \begin{gathered}
 N_e=[\,\sigma_{e,\ell_1}m_{\ell_1}^{\mathsf T}\ \cdots\
       \sigma_{e,\ell_{21}}m_{\ell_{21}}^{\mathsf T}\,],\qquad \Delta_e=\det N_e,\\
 d_e=c^{\mathsf T}-\sum_{i=1}^{40}\sigma_{e,i}m_i^{\mathsf T}.
 \end{gathered}
\]
Let $N_e[r\leftarrow d_e]$ mean replace column $r$ of $N_e$ by $d_e$, and put
\[
 \lambda_{e,i}=\begin{cases}
 1+\det(N_e[r\leftarrow d_e])/\Delta_e,&i=\ell_r\ (1\le r\le21),\\
 1,&i\notin B_e.
 \end{cases}
\]
Evaluating these integer determinants exactly gives
\[
 \Delta_e\ne0\quad(0\le e\le40),\qquad
 \frac{\det(N_e[r\leftarrow d_e])}{\Delta_e}>8
 \quad(0\le e\le40,\ 1\le r\le21).
\]
Thus all weights are at least $1$ (basis weights exceed $9$), and Cramer's rule gives
$c=\sum_{i=1}^{40}\lambda_{e,i}\sigma_{e,i}m_i$ for every $e$.
Also $\Delta_0\ne0$ gives $\operatorname{rank}M_0=21$. If $z\ne0$ has all
signed slacks $\sigma_{e,i}m_i z\ge0$, full column rank makes at least one
slack positive, so $c(z)>0$. Every datum needed for these finite checks is
printed in this appendix. An exact checker of these formulas is included
in \cite{supplement}.
% END EXPLICIT CERTIFICATE

For example, the first row is
$m_1=u_1+100u_2+21u_5-7u_6+u_{17}+u_{18}$.
The supplement's exact checker reconstructs every entry from these
recipes, compares them with the rational inequality data, and verifies
the displayed weight identities. The positive functional used in Lean
differs from the one printed here.

\section{Completeness of the vertex certificate}\label{app:cover}

The following local certificate criterion explains the completeness fact
used in Section~3. The large finite list is in \cite{supplement};
completeness is a consequence, not a premise supplied by the list.
The labeled pair, potential, path and local source data are supplied
only in the supplement and checked there; their large tables are not
printed here. The pair has the complementary-zero-set form used in
\cite[proof of Proposition~8]{adkz}.

\begin{lemma}\label{lem:cover}
Let $\mathcal Q=\{x\in\R^d:Lx\ge\zeta\}$ be full-dimensional.
Suppose a nonempty finite set $\mathcal V\subset\mathcal Q$ is given, each $v\in \mathcal V$
having exactly $d$ tight rows, whose matrix $L_v$ is invertible.
For each of these rows $j$, suppose a point $v_j\in \mathcal V$ satisfies
\[
 L_v(v_j-v)=g_j\mathbf e_j,\qquad g_j>0,
\]
where $\mathbf e_j$ is the $j$th coordinate vector. Then $\mathcal V$ is the complete
vertex set, $\mathcal Q$ is simple, and its neighbors at $v$ are precisely
the points $v_j$.
\end{lemma}
\begin{proof}
For a linear objective $f$, choose $v\in \mathcal V$ maximizing $f$ on this finite
set. For every $x\in\mathcal Q$, the vector $\kappa=L_v(x-v)$ has
nonnegative coordinates and
\[
 x-v=\sum_{j=1}^d\frac{\kappa_j}{g_j}(v_j-v).
\]
Since $f(v_j-v)\le0$, this proves $f(x)\le f(v)$: the list contains a
global maximizer for every linear objective.

Let $x$ be any vertex and choose $d$ independent tight rows $I$ there.
The objective $f(y)=-\sum_{i\in I}L_i y$ is maximized precisely where
all these rows are tight, because their slacks are nonnegative.
Its unique feasible maximizer is $x$. A listed maximizer must therefore
equal $x$. Conversely each listed point is a vertex by its nonsingular
tight system, so the list is complete and every vertex has exactly $d$
tight rows. Each such row defines a distinct facet locally. Releasing one
row while retaining the other $d-1$ defines an edge. The distinct vertices
$v,v_j$ lie on that one-dimensional face and hence are its two endpoints.
Every edge at a simple vertex arises this way.
\end{proof}

Here are the exact formal interfaces for these inputs; all names have
prefix \nolinkurl{ADKZPricing}.
For the literal matrix $M_0$, \nolinkurl{Generated.source_valid} proves
\nolinkurl{IntegerSourceValid} for the supplied 36,425 rays, their
20-element bases and their neighbors. Its field hypotheses are checked
finite arithmetic: positive normalizing coordinates, feasible rays,
injective and exact active labels, and a neighbor retaining every basis
row except the released one. The resulting
\nolinkurl{Generated.sourceCertificate} is a
\nolinkurl{SourceVertexCertificate}.
For this certificate, the theorems
\nolinkurl{SourceVertexCertificate.covers_active_labels},
\nolinkurl{SourceVertexCertificate.active_count} and
\nolinkurl{SourceVertexCertificate.basis_nonsingular} give
\[
 \operatorname{Vert}(P)=\{w_v:v\in\mathcal V\},\qquad
 |I(w_v)|=20,\qquad
 \operatorname{rank}(M_0)_{I(w_v)}=20,
\]
where $\mathcal V$ is the supplied list, $w_v$ its ray normalized by
$c(w_v)=1$, and $I(w_v)$ its tight-row set. In particular, $P$ is simple.
The normalization changes only the scale of each ray, preserving vertices,
facets and their incidence; Lemma~\ref{lem:cover} is the corresponding
written local-to-global argument.

The theorem \nolinkurl{msw_integral_real_vertices} verifies the endpoints
with tight rows 1--20 and 21--40.
The theorems \nolinkurl{Generated.source_local_potential} and
\nolinkurl{Generated.source_specified_gap} verify the local potential
inequalities and endpoint gap at least 21.
Applying \nolinkurl{SourceVertexCertificate.potential_edge} and
\nolinkurl{SourceVertexCertificate.specified_distance_lower_bound}
therefore gives a 1-Lipschitz potential on every geometric edge and
$d_P(u,v)\ge21$, with no assumed completeness or distance premise.
The separate exact checkers \texttt{check\_fraction.py} and
\texttt{check\_rational.py} verify the supplied 21-edge path, giving
$d_P(u,v)=21$; this upper bound is not a Lean distance-equality theorem.
These source facts and Appendix~\ref{app:certificate}'s identities are
the complete computational input to the written construction.
No enumeration of the final 80-element cocircuit graph is used.

\section{Earlier farthest-pair questions}

The diameter conjecture is related to earlier questions about which pairs
of cocircuits are farthest apart; we distinguish those questions here.
For uniform oriented matroids, Babson, Finschi and Fukuda prove antipodal distance $n-r+2$
\cite[Lemma~2]{bff} and equality exactly for antipodal pairs in uniform
ranks two and three \cite[Lemma~6]{bff}. ADKZ answer Finschi's Open Problem~2
and refute the equality half of \cite[Question~4.2]{felsner}
with their nonantipodal example \cite[Proposition~8 and preceding paragraph]{adkz}.
Their proof already uses a pair at distance 21 in the source polytope;
here an explicit path gives equality for the source matrix specified in Appendix~\ref{app:certificate}.
Finschi's Open Problem~4 asks whether a uniform cocircuit graph can
have an antipodal pair whose distance differs from the diameter, or a
nonantipodal pair attaining the diameter \cite[p.~90]{finschi}; this is
the farthest-pairs question in \cite[final remarks, pp.~13--14]{bff}.
ADKZ's lower-bound argument answers it affirmatively once the source
example is made uniform, as in Lemmas~\ref{lem:facetcross}--\ref{lem:factor}:
if the diameter is 21, the nonantipodal pair attains it; if larger,
antipodal pairs are not farthest.
Here every antipodal pair is shorter than some nonantipodal pair.
This refutes the diameter half of Question~4.2 and ADKZ's Conjecture~1.

A \emph{crabbed path} from $X$ to $Y$ has
$Z^+\subseteq X^+\cup Y^+$ and $Z^-\subseteq X^-\cup Y^-$ at every
vertex $Z$, where $Z^+=\{i:Z_i=1\}$ and $Z^-=\{i:Z_i=-1\}$.
ADKZ exhibit a realizable rank-four example on nine elements in which a
shortest path is shorter than every crabbed path \cite[Theorem~9]{adkz}.

Adjacent work concerns e-positive circuit graphs of oriented matroid
polyhedra and embracing exchange distance
(B\'erczi, Knauer, Montejano and N\'ador \cite[Corollary~1.4 and Definition~2.1]{embracing}),
smaller non-Hirsch polytopes (\'Swirszcz and collaborators \cite[\S4.1]{hopper}),
and polyhedron certification (Allamigeon, Id-Sahra and Strub
\cite[abstract and \S1]{cert26}).

\end{appendices}

\begin{thebibliography}{99}
\footnotesize\raggedright
\setlength{\baselineskip}{10pt}
\setlength{\itemsep}{0pt}
\bibitem{bff}
E.~Babson, L.~Finschi and K.~Fukuda,
\emph{Cocircuit graphs and efficient orientation reconstruction in oriented
matroids}, European Journal of Combinatorics \textbf{22} (2001), 587--600.
\href{https://doi.org/10.1006/eujc.2001.0481}{doi:10.1006/eujc.2001.0481}.
Statement locators refer to the \href{https://finschi.com/math/publ/2001-05-04_Babson-Finschi-Fukuda_Cocircuit-Graphs-Orientation-Reconstruction.pdf}{May~4, 2001 author version}.
\bibitem{finschi}
L.~Finschi, \emph{A graph theoretical approach for reconstruction and
generation of oriented matroids}, doctoral thesis, ETH Z\"urich, 2001.
\href{https://doi.org/10.3929/ethz-a-004255224}{doi:10.3929/ethz-a-004255224}.
\href{https://finschi.com/math/publ/2001-08-31_Finschi_A-Graph-Theoretical-Approach-for-Reconstruction-and-Generation-of-Oriented-Matroids.pdf}{Author copy}.
\bibitem{felsner}
S.~Felsner, R.~G\'omez, K.~Knauer, J.~J.~Montellano-Ballesteros and R.~Strausz,
\emph{Cubic time recognition of cocircuit graphs of uniform oriented matroids},
European Journal of Combinatorics \textbf{32} (2011), 60--66.
\href{https://doi.org/10.1016/j.ejc.2010.07.012}{doi:10.1016/j.ejc.2010.07.012}.
\href{https://arxiv.org/abs/1001.3534v1}{arXiv:1001.3534v1} (2010;
statement numbering).
\bibitem{adkz}
I.~Adler, J.~A.~De~Loera, S.~Klee and Z.~Zhang,
\emph{Diameters of cocircuit graphs of oriented matroids: an update},
Electronic Journal of Combinatorics \textbf{28} (2021), no.~4, P4.40.
\href{https://doi.org/10.37236/9653}{doi:10.37236/9653}.
\bibitem{survey}
G.~M.~Ziegler, L.~Anderson and K.~Knauer,
\emph{Oriented matroids today}, Electronic Journal of Combinatorics,
Dynamic Survey DS4, version~4, April~16, 2024.
\href{https://www.combinatorics.org/files/Surveys/ds4/ds4v4-2024.pdf}{Version~4}.
\bibitem{msw}
B.~Matschke, F.~Santos and C.~Weibel,
\emph{The width of five-dimensional prismatoids}, Proceedings of the London
Mathematical Society (3) \textbf{110} (2015), 647--672.
\href{https://doi.org/10.1112/plms/pdu064}{doi:10.1112/plms/pdu064}.
\href{https://arxiv.org/abs/1202.4701v2}{arXiv:1202.4701v2}
(statement numbering).
\bibitem{santos}
F.~Santos, \emph{A counterexample to the Hirsch conjecture},
Annals of Mathematics \textbf{176} (2012), 383--412.
\href{https://doi.org/10.4007/annals.2012.176.1.7}{doi:10.4007/annals.2012.176.1.7}.
Statement numbering refers to the Annals version.
\bibitem{acs}
X.~Allamigeon, Q.~Canu and P.-Y.~Strub,
\emph{A Formal Disproof of the Hirsch Conjecture}, Proceedings of the 12th ACM
SIGPLAN International Conference on Certified Programs and Proofs (2023), 17--29.
\href{https://doi.org/10.1145/3573105.3575678}{doi:10.1145/3573105.3575678}.
\href{https://arxiv.org/abs/2301.04060v1}{arXiv:2301.04060v1}
(statement and page numbering).
\bibitem{embracing}
K.~B\'erczi, K.~Knauer, L.~P.~Montejano and B.~N\'ador,
\emph{Embracing exchange sequences and oriented matroid polyhedron diameter},
\href{https://arxiv.org/abs/2606.19573v2}{arXiv:2606.19573v2} (2026).
\bibitem{hopper}
G.~\'Swirszcz, A.~Z.~Wagner, G.~Williamson, S.~Blackwell, B.~Georgiev,
A.~Davies, A.~Eslami, S.~Racani\`ere, T.~Weber and P.~Kohli,
\emph{Advancing Geometry with AI: Multi-agent Generation of Polytopes},
\href{https://arxiv.org/abs/2502.05199v1}{arXiv:2502.05199v1} (2025).
\bibitem{cert26}
X.~Allamigeon, Y.~Id-Sahra and P.-Y.~Strub,
\emph{Formally Certifying the Vertex Set of a Polyhedron Faster than Informal
Enumeration},
\href{https://arxiv.org/abs/2610.11913v1}{arXiv:2610.11913v1} (2026).
\bibitem{supplement}
P.~Graham, \emph{Supplementary material for A counterexample to the cocircuit
diameter conjecture}, version~1 (2026).
\href{https://github.com/6/clank-lab/tree/cocircuit-diameter-counterexample-v1/results/cocircuit-diameter-counterexample}{Supplementary data and Lean formalization}.
\bibitem{klee}
S.~Klee, \emph{Diameter bounds on cocircuit graphs of oriented matroids},
Oberwolfach Report 39/2019, pp.~2424--2426 (joint work with I.~Adler,
J.~A.~De~Loera and Z.~Zhang).
\href{https://ems.press/content/serial-article-files/46818}{Report}.
\end{thebibliography}
\end{document}
